\documentclass[12pt]{book}
\usepackage{precalc1}
\setcounter{chapter}{9}
\begin{document}

\chapter{Matrices and Gaussian Elimination}

Solving a system by elimination, the variable names $x,y,z$ are carried along
unchanged through every step; only the coefficients move. A matrix strips away the
names and keeps the numbers in a grid, so the bookkeeping is done on coefficients
alone. The elimination procedure then becomes a fixed sequence of row operations on
that grid, and the method, applied systematically, is Gaussian elimination.

\section{The Augmented Matrix}

A linear system is summarized by its coefficients and right-hand sides, arranged in a
rectangular array. Each row is an equation; each column before the bar is a variable;
the column after the bar holds the constants.

\begin{definition}{Augmented matrix}
The \emph{augmented matrix} of a linear system records each equation's coefficients
in a row, aligned by variable in columns, with the constant terms placed after a
vertical bar. The system
\[
\begin{aligned}
2x+y-z&=1\\ x-y+2z&=3\\ 3x+2y+z&=10
\end{aligned}
\qquad\text{becomes}\qquad
\left[\begin{array}{ccc|c}
2&1&-1&1\\ 1&-1&2&3\\ 3&2&1&10
\end{array}\right].
\]
\end{definition}

\begin{keyidea}
The augmented matrix is pure bookkeeping: the same system with the variable names
removed. Every operation legal on equations corresponds to an operation on rows, and
nothing is lost, because the column positions remember which variable each number
belongs to.
\end{keyidea}

\section{Row Operations}

The three reversible moves on equations, from the previous chapter, become the three
elementary row operations. Each leaves the solution set unchanged.

\begin{keyidea}
\textbf{Elementary row operations.}
\begin{enumerate}\setlength{\itemsep}{2pt}
  \item \emph{Swap} two rows (reorder the equations).
  \item \emph{Scale} a row by a nonzero constant (multiply an equation through).
  \item \emph{Replace} a row by itself plus a multiple of another row (add a multiple
        of one equation to another).
\end{enumerate}
\end{keyidea}

The goal of these operations is \emph{row echelon form}: a staircase pattern in which
each row's first nonzero entry (its leading entry) sits to the right of the one
above, and entries below each leading position are zero. A matrix in this form is the
matrix version of a triangular system, ready for back-substitution.

\begin{visual}
Row echelon form is a staircase of zeros below the diagonal. The leading entry of
each row steps rightward as you move down, and everything beneath a leading entry is
cleared to zero. Reading the bottom row gives one variable; reading upward recovers
the rest.
\end{visual}

\section{Gaussian Elimination}

Gaussian elimination is the disciplined use of row operations to reach echelon form,
working one column at a time from left to right: use the leading entry of a row to
clear all entries below it, then move to the next column. Once the staircase is
complete, back-substitution finishes the solve.

\begin{example}[a full elimination]
Solve
\[
\left[\begin{array}{ccc|c}
1&1&1&6\\ 2&-1&1&3\\ 1&2&-1&2
\end{array}\right].
\]
Clear the first column below the leading $1$. Replace row $2$ by row $2$ minus twice
row $1$, and row $3$ by row $3$ minus row $1$:
\[
\left[\begin{array}{ccc|c}
1&1&1&6\\ 0&-3&-1&-9\\ 0&1&-2&-4
\end{array}\right].
\]
Clear the second column below the new leading entry. Swap rows $2$ and $3$ for a
tidy leading $1$, then replace row $3$ by row $3$ plus three times row $2$:
\[
\left[\begin{array}{ccc|c}
1&1&1&6\\ 0&1&-2&-4\\ 0&-3&-1&-9
\end{array}\right]
\rightarrow
\left[\begin{array}{ccc|c}
1&1&1&6\\ 0&1&-2&-4\\ 0&0&-7&-21
\end{array}\right].
\]
The matrix is in echelon form. The bottom row reads $-7z=-21$, so $z=3$. The middle
row reads $y-2z=-4$, so $y=2$. The top row reads $x+y+z=6$, so $x=1$. The solution is
$(1,2,3)$, matching the same system solved by hand in the previous chapter.
\end{example}

\begin{keyidea}
Gaussian elimination reaches echelon form column by column: pick a leading entry,
clear everything below it, advance. The procedure is mechanical and always
terminates, which is why it scales to systems far larger than hand elimination can
manage.
\end{keyidea}

\section{Reading the Outcome from the Matrix}

The three outcomes appear in the echelon form as clearly as they did in the geometry.
A full staircase with a leading entry in every variable column gives a unique
solution. A row of zeros in the coefficients with a nonzero constant, the pattern
$[\,0\ 0\ 0\mid c\,]$ with $c\neq0$, reads $0=c$, a contradiction, so no solution. A
row of all zeros, $[\,0\ 0\ 0\mid 0\,]$, reads $0=0$, carrying no information and
leaving a free variable, so infinitely many.

\begin{keyidea}
\textbf{Reading echelon form.}
\begin{itemize}\setlength{\itemsep}{2pt}
  \item A leading entry in every variable column: \emph{unique solution}.
  \item A row $[\,0\ \cdots\ 0\mid c\,]$ with $c\neq0$: \emph{no solution}
        (the false statement $0=c$).
  \item A row of all zeros and a column without a leading entry: \emph{infinitely
        many}, with a free variable.
\end{itemize}
\end{keyidea}

\begin{example}[a system with no solution]
Eliminating
\[
\left[\begin{array}{cc|c}
1&-1&2\\ 2&-2&7
\end{array}\right]
\]
by replacing row $2$ with row $2$ minus twice row $1$ gives
\[
\left[\begin{array}{cc|c}
1&-1&2\\ 0&0&3
\end{array}\right].
\]
The bottom row reads $0=3$, false, so the system has no solution; the two lines are
parallel.
\end{example}

\begin{example}[a system with infinitely many]
Eliminating
\[
\left[\begin{array}{cc|c}
1&2&3\\ 2&4&6
\end{array}\right]
\]
gives a bottom row of all zeros, $0=0$. The variable $y$ is free; setting $y=t$, the
top row gives $x=3-2t$. The solution set is $(3-2t,\,t)$, the whole line, as before.
\end{example}

\section{Applications}

\subsection{Currents in a circuit}

Kirchhoff's two laws turn a circuit into a linear system. The \emph{current law}
states that the current into any junction equals the current out; the \emph{voltage
law} states that around any closed loop, the sum of voltage rises from batteries
equals the sum of drops across resistors, each drop being resistance times current
($V=IR$). Applying both laws produces one equation per junction and one per loop.

\begin{center}
\begin{tikzpicture}[scale=1.0,line width=0.9pt]
  % outer rectangle
  \draw (0,0) rectangle (5,2.4);
  % middle branch
  \draw (2.5,0)--(2.5,2.4);
  % battery E1 left side
  \node[left] at (0,1.2) {};
  \draw (-0.05,1.05)--(-0.05,1.35); \draw (0.05,0.95)--(0.05,1.45);
  \node[left] at (-0.1,1.2) {\scriptsize $E_1$};
  % battery E2 right side
  \draw (4.95,1.05)--(4.95,1.35); \draw (5.05,0.95)--(5.05,1.45);
  \node[right] at (5.1,1.2) {\scriptsize $E_2$};
  % resistors (zigzag as small boxes)
  \draw[fill=white] (1.0,2.3) rectangle (1.7,2.5); \node at (1.35,2.72) {\scriptsize $R_1$};
  \draw[fill=white] (3.3,2.3) rectangle (4.0,2.5); \node at (3.65,2.72) {\scriptsize $R_3$};
  \draw[fill=white] (2.4,1.0) rectangle (2.6,1.7); \node at (2.95,1.35) {\scriptsize $R_m$};
  % current labels
  \node[pcblue] at (0.5,1.2) {\scriptsize $I_1$};
  \node[pcblue] at (2.5,0.55) {\scriptsize $I_2$};
  \node[pcblue] at (4.5,1.2) {\scriptsize $I_3$};
  % top node
  \fill (2.5,2.4) circle (1.6pt); \node[above right] at (2.5,2.4) {\scriptsize node};
\end{tikzpicture}
\end{center}

\begin{example}[a two-loop circuit]
A circuit has a left branch carrying current $I_1$ through a battery $E_1=12$ V and
resistor $R_1=2\,\Omega$, a middle branch carrying $I_2$ through $R_m=3\,\Omega$, and
a right branch carrying $I_3$ through $R_3=4\,\Omega$ and battery $E_2=10$ V. The
current law at the top node and the voltage law around the two loops give
\[
\begin{aligned}
I_1-I_2-I_3&=0 &&\text{(junction)},\\
2I_1+3I_2&=12 &&\text{(left loop)},\\
3I_2+4I_3&=10 &&\text{(right loop)}.
\end{aligned}
\]
The augmented matrix and its elimination:
\[
\left[\begin{array}{ccc|c}
1&-1&-1&0\\ 2&3&0&12\\ 0&3&4&10
\end{array}\right]
\xrightarrow{\,R_2-2R_1\,}
\left[\begin{array}{ccc|c}
1&-1&-1&0\\ 0&5&2&12\\ 0&3&4&10
\end{array}\right]
\xrightarrow{\,5R_3-3R_2\,}
\left[\begin{array}{ccc|c}
1&-1&-1&0\\ 0&5&2&12\\ 0&0&14&14
\end{array}\right].
\]
The bottom row gives $14I_3=14$, so $I_3=1$ A. Back-substituting, $5I_2+2=12$ gives
$I_2=2$ A, and $I_1-2-1=0$ gives $I_1=3$ A. The currents are $I_1=3$, $I_2=2$,
$I_3=1$ amperes; the junction law checks, since $3=2+1$.
\end{example}

\subsection{Balancing a chemical equation}

\begin{example}[combustion of propane]
Balance $a\,\mathrm{C_3H_8}+b\,\mathrm{O_2}\to c\,\mathrm{CO_2}+d\,\mathrm{H_2O}$.
Equal atom counts on each side give one equation per element:
\[
\begin{aligned}
\text{C:}&\quad 3a=c,\\
\text{H:}&\quad 8a=2d,\\
\text{O:}&\quad 2b=2c+d.
\end{aligned}
\]
This is a homogeneous system in $a,b,c,d$. Take $a=t$ as the free parameter. The
carbon equation gives $c=3t$; the hydrogen equation gives $d=4t$; substituting into
the oxygen equation gives $2b=2(3t)+4t=10t$, so $b=5t$. The smallest whole numbers
come from $t=1$: $a=1$, $b=5$, $c=3$, $d=4$, yielding
\[
\mathrm{C_3H_8}+5\,\mathrm{O_2}\to3\,\mathrm{CO_2}+4\,\mathrm{H_2O}.
\]
The free parameter is exactly the freedom to scale a balanced equation up or down.
\end{example}

\subsection{Flow through a road network}

\begin{example}[a one-way traffic grid]
Four one-way streets meet at four intersections $A,B,C,D$, with the unknown segment
flows $x_1,x_2,x_3,x_4$ in cars per hour. Conservation at each intersection (flow in
equals flow out), with the external flows shown by the data, gives
\[
\begin{aligned}
A:&\quad x_1+x_2=800,\\
B:&\quad x_2=x_3+200,\\
C:&\quad x_3+x_4=600,\\
D:&\quad x_1=x_4.
\end{aligned}
\]
Eliminating leaves one row of zeros: the four conservation equations are not
independent, because total flow in must equal total flow out, so one equation follows
from the others. The system is underdetermined, with one free parameter. Taking
$x_4=t$, back-substitution gives
\[
x_1=t,\qquad x_2=800-t,\qquad x_3=600-t,\qquad x_4=t.
\]
The flows are not uniquely fixed; the network admits a family of balanced states. The
physical requirement that each flow be nonnegative forces $0\le t\le600$, the range of
admissible traffic patterns. A measurement of any one segment pins down all four.
\end{example}

\section{Exercises}

\begin{exercises}
\item Write the augmented matrix of
$\begin{aligned}3x-y+2z&=5\\ x+4y-z&=0\\ 2x+y+z&=7\end{aligned}$, labeling which
column belongs to which variable.

\item State the three elementary row operations and the equation operation each one
represents.

\item Apply one row operation to clear the first column below the top entry of
$\left[\begin{array}{cc|c}1&3&5\\ 2&1&4\end{array}\right]$, and write the result.

\item Solve by Gaussian elimination:
$\left[\begin{array}{cc|c}1&2&5\\ 3&1&5\end{array}\right]$, and state the solution as
a point.

\item Carry
$\left[\begin{array}{ccc|c}1&1&1&4\\ 2&1&-1&3\\ 1&-2&3&11\end{array}\right]$
to echelon form and solve by back-substitution; check your answer.

\item Eliminate
$\left[\begin{array}{cc|c}1&-2&3\\ -3&6&5\end{array}\right]$
and state the outcome, identifying the row that reveals it.

\item Eliminate
$\left[\begin{array}{cc|c}2&4&6\\ 3&6&9\end{array}\right]$,
identify the free variable, and write the solution set with a parameter.

\item Explain how each of the three outcomes, unique, none, and infinitely many,
appears in the echelon form of the augmented matrix.

\item A system's echelon form has a bottom row $[\,0\ 0\ 0\mid 5\,]$. State the
outcome and explain the reasoning in one sentence.

\item Balance $a\,\mathrm{Fe}+b\,\mathrm{O_2}\to c\,\mathrm{Fe_2O_3}$ by writing the
atom-count equations as a linear system and choosing the smallest whole-number
solution.

\item A two-loop circuit gives the system
$\begin{aligned}I_1-I_2-I_3&=0\\ I_1+2I_2&=10\\ 2I_2+3I_3&=9\end{aligned}$
for its branch currents. Solve by Gaussian elimination, and verify the junction law
$I_1=I_2+I_3$.

\item Explain why Gaussian elimination on a matrix and hand elimination on the
equations produce the same solution, and why the matrix form is preferred for large
systems.
\end{exercises}

\end{document}
